% Part: first-order-logic
% Chapter: beyond
% Section: higher-order-logic

\documentclass[../../../include/open-logic-section]{subfiles}

\begin{document}

\olfileid{fol}{byd}{hol}

\olsection{ఉన్నత-స్థాయి తర్కం}

మొదటిస్థాయి తర్కం నుంచి ద్వితీయ-స్థాయి తర్కానికి వెళ్లడం వల్ల,
ఔపచారిక భాషలోనే మొదటిస్థాయి \tetoken{వ్యక్తి క్షేత్రంలోని}{!!{domain}} వస్తువుల సమితుల గురించి
మాట్లాడగలిగాం. అక్కడే ఎందుకు ఆగాలి? ఉదాహరణకు, వస్తువుల సమితుల
సమితులతో, బహుశా వస్తువులనూ వస్తువుల సమితులనూ రెండింటినీ కలిగిన
సమితులతో కూడా వ్యవహరించడానికి తృతీయ-స్థాయి తర్కం వీలు కల్పించాలి.
అటువంటి వస్తువుల సమితుల గురించి మాట్లాడటానికి చతుర్థ-స్థాయి తర్కం
వీలు కల్పిస్తుంది. మీరు ఊహించినట్లే, ఈ ఆలోచనను ఎన్ని సార్లైనా
పునరావృతం చేయవచ్చు.

ఆచరణలో, ఉన్నత-స్థాయి తర్కాన్ని సంబంధాలకు బదులు ప్రమేయాల ఆధారంగా
తరచుగా \tetoken{సూత్రంగా}{!!{formula}} రూపొందిస్తారు. (సహజ గుర్తింపులను మినహాయిస్తే ఈ
తేడా అనావశ్యకం.) కొన్ని మౌలిక ``రకాలు'' $A$, $B$, $C$,~\dots
(ఇప్పటి నుంచి వీటిని ``టైపులు'' అంటాం) ఇచ్చినప్పుడు, కింది నిబంధనతో
కొత్తవాటిని సృష్టించవచ్చు:
\begin{quote}
$\sigma$, $\tau$లు పరిమిత టైపులు అయితే, $\sigma \to \tau$ కూడా పరిమిత
టైపే.
\end{quote}
టైపులను మన \tetoken{వ్యక్తి క్షేత్రంలో}{!!{domain}} కావలసిన వస్తువులను వర్గీకరించే
వాక్యనిర్మాణాత్మక ``గుర్తులుగా'' భావించండి; టైపు~$\sigma$ వస్తువులను
టైపు~$\tau$ వస్తువులకు పంపే ప్రమేయాలైన వస్తువులను
$\sigma \to \tau$ వివరిస్తుంది. ఉదాహరణకు, ``సత్యం'', ``అసత్యం'' అనే
సత్యమూల్యాల టైపు~$\Omega$, సహజ సంఖ్యల టైపు~$\Nat$ కావాలనుకోవచ్చు.
అప్పుడు $\Nat \to \Omega$ టైపు వస్తువులను ఏకస్థాన సంబంధాలుగా లేదా
$\Nat$ ఉపసమితులుగా భావించవచ్చు; $\Nat \to \Nat$ టైపు వస్తువులు సహజ
సంఖ్యల నుంచి సహజ సంఖ్యలకు ప్రమేయాలు; $(\Nat \to \Nat) \to \Nat$
టైపు వస్తువులు ``ప్రమేయకాలు'', అంటే ప్రమేయాలను సంఖ్యలకు పంపే
ఉన్నత-టైపు ప్రమేయాలు.

ద్వితీయ-స్థాయి తర్కంలోలాగే, మనం పరిగణించాలనుకునే ప్రతి వస్తు టైపుకూ
ఒక రకం ఉండే బహు-రక తర్క రూపంగా ఉన్నత-స్థాయి తర్కాన్ని భావించవచ్చు.
అయితే ఉన్నత-టైపు తర్కపు వాక్యనిర్మాణాన్ని మొదటి నుంచి నిర్వచించడం
సాధారణంగా మరింత స్పష్టం. ఉదాహరణకు, పరిమిత టైపుల సమితిని కింది విధంగా
ఆగమన పద్ధతిలో నిర్వచించవచ్చు:
\begin{enumerate}
\item $\Nat$ ఒక పరిమిత టైపు.
\item $\sigma$, $\tau$లు పరిమిత టైపులు అయితే, $\sigma
  \to \tau$ కూడా పరిమిత టైపే.
\item $\sigma$, $\tau$లు పరిమిత టైపులు అయితే, $\sigma \times
  \tau$ కూడా పరిమిత టైపే.
\end{enumerate}
అంతఃప్రజ్ఞాత్మకంగా, $\Nat$ సహజ సంఖ్యల టైపును,
$\sigma \to \tau$ అనేది $\sigma$ నుంచి $\tau$కు ప్రమేయాల టైపును,
$\sigma \times \tau$ అనేది $\sigma$ నుంచి ఒకటి, $\tau$ నుంచి ఒకటి
చొప్పున తీసుకున్న వస్తు జతల టైపును సూచిస్తాయి. తరువాత పదాల సమితిని
కింది విధంగా ఆగమన పద్ధతిలో నిర్వచించవచ్చు:
\begin{enumerate}
\item ప్రతి టైపు~$\sigma$కు, టైపు~$\sigma$ గల $x$, $y$, $z$, \dots
  చరాల నిల్వ ఒకటి ఉంటుంది.
\item $\Obj 0$ అనేది టైపు~$\Nat$ గల పదం.
\item $\Obj S$ (ఉత్తరాధికారి) అనేది టైపు~$\Nat \to \Nat$ గల పదం.
\item $s$ టైపు~$\sigma$ గల పదం, $t$ టైపు
  $\Nat \to (\sigma \to \sigma)$ గల పదం అయితే, $\Obj{R}_{st}$ టైపు
  $\Nat \to \sigma$ గల పదం.
\item $s$ టైపు~$\tau \to \sigma$ గల పదం, $t$ టైపు~$\tau$ గల పదం
  అయితే, $s(t)$ టైపు~$\sigma$ గల పదం.
\item $s$ టైపు~$\sigma$ గల పదం, $x$ టైపు~$\tau$ గల చరం అయితే,
  $\lambd[x][s]$ టైపు~$\tau \to \sigma$ గల పదం.
\item $s$ టైపు~$\sigma$ గల పదం, $t$ టైపు~$\tau$ గల పదం అయితే,
  $\tuple{s, t}$ టైపు~$\sigma \times \tau$ గల పదం.
\item $s$ టైపు~$\sigma \times \tau$ గల పదం అయితే, $p_1(s)$
  టైపు~$\sigma$ గల పదం, $p_2(s)$ టైపు~$\tau$ గల పదం.
\end{enumerate}
అంతఃప్రజ్ఞాత్మకంగా, కింది పునరావర్తనంతో నిర్వచించిన ప్రమేయాన్ని
$\Obj{R}_{st}$ సూచిస్తుంది:
\begin{align*}
\Obj{R}_{st}(0) & = s \\
\Obj{R}_{st}(x+1) & = t(x, R_{st}(x)),
\end{align*}
$\tuple{s, t}$ అనేది మొదటి అంగం~$s$, రెండవ అంగం~$t$ గల జతను
సూచిస్తుంది; $p_1(s)$, $p_2(s)$లు $s$ యొక్క మొదటి, రెండవ మూలకాలను
(``ప్రక్షేపాలను'') సూచిస్తాయి. చివరగా, టైపు~$\tau$ గల ఏ $x$కైనా
కిందిదిగా నిర్వచించిన ప్రమేయం~$f$ను $\lambd[x][s]$ సూచిస్తుంది:
\[
f(x) = s
\]
అందువల్ల అంశం (6), పదాలను వాడి ప్రమేయాలను నిర్వచించడానికి వీలు
కల్పించే ఒక ధర్మసంగ్రహ రూపాన్ని ఇస్తుంది.
\sourcecorrection{OLTEFOLBYD-004}{ఆంగ్ల మూలం లాంబ్డా పదంలోని బద్ధ చరం $x$ టైపు~$\sigma$ గలదని చెప్పింది; ముందున్న అంశం (6)లో $x$ టైపు~$\tau$, ఫలిత పదం $s$ టైపు~$\sigma$ అని నిర్దేశించినందున ఇక్కడ $\tau$గా సరిచేశాం.}
ఒకే టైపు గల పదాల మధ్య \tetoken{తాదాత్మ్య విధేయం}{!!{identity}} ప్రకటనలు $\eq[s][t]$, సాధారణ
ప్రతిజ్ఞావాక్య సంయోజకాలు, ఉన్నత-టైపు పరిమాణీకరణ నుంచి
\tetoken{సూత్రాలను}{!!^{formula}s} నిర్మిస్తాం. పైన నిర్వచించిన పదాలను నియంత్రించే
మౌలిక సమీకరణాలను, పరిమాణకాలు, \tetoken{తాదాత్మ్య విధేయం}{!!{identity}} గల సాధారణ తర్క
నియమాలను కలిపి వ్యవస్థ స్వీకృతాలుగా తీసుకోవచ్చు.

పరిమిత టైపు వ్యవస్థకు సత్యమూల్యాల టైపు~$\Omega$ను చేరిస్తే, దాని
వాడకాన్ని నియంత్రించే స్వీకృతాలను కూడా చేర్చాలి. నిజానికి చాకచక్యంగా
వ్యవహరిస్తే, సంక్లిష్ట \tetoken{సూత్రాలను}{!!{formula}s} పూర్తిగా తొలగించి, వాటి స్థానంలో
టైపు~$\Omega$ గల పదాలను ఉంచవచ్చు!{} అప్పుడు నిరూపణ వ్యవస్థను దానికి
అనుగుణంగా మార్చవచ్చు. ఫలితంగా, 1930లలో అలోంజో చర్చ్ ప్రతిపాదించిన
\emph{సరళ టైపుల సిద్ధాంతం} సారాంశంలో లభిస్తుంది.

ద్వితీయ-స్థాయి తర్కంలోలాగే, మనం వాడాలనుకునే ఉన్నత-టైపు
అర్థవిచారానికి వేర్వేరు రూపాలు ఉన్నాయి. పూర్ణ రూపంలో,
టైపు~$\sigma$ వస్తువుల నుంచి టైపు~$\tau$ వస్తువులకు గల
\emph{అన్ని} ప్రమేయాల సమితిపై టైపు~$\sigma \to \tau$ చరాలు
విస్తరిస్తాయి. ఊహించగలిగినట్లే, ఒక సంపూర్ణ, ప్రభావవంతమైన
\tetoken{వ్యుత్పత్తి}{!!{derivation}} వ్యవస్థను అనుమతించలేనంత బలమైన అర్థవిచారం ఇది. అయితే,
ప్రతి టైపు~$\tau$కు $T_\tau$ అనే మూలకాల సమితులు, ప్రయోగం,
ప్రక్షేపం మొదలైన వాటికి తగిన క్రియలు కలసి \tetoken{ఒక నిర్మాణం}{!!a{structure}} అయ్యే
బలహీన అర్థవిచారాన్ని పరిగణించవచ్చు. వివరాలను సరిగ్గా అమలు చేస్తే,
పైన వివరించిన రకాల \tetoken{వ్యుత్పత్తి}{!!{derivation}} వ్యవస్థలకు సంపూర్ణతా సిద్ధాంతాలను
పొందవచ్చు.

ఉన్నత-టైపు తర్కం ఆకర్షణీయమైనది; ఎందుకంటే గణితంలో గణనీయమైన భాగాన్ని
సహజంగా ఇమిడ్చగల చట్రాన్ని అది అందిస్తుంది: $\Nat$తో మొదలుపెట్టి
వాస్తవ సంఖ్యలు, అవిచ్ఛిన్న ప్రమేయాలు మొదలైనవాటిని నిర్వచించవచ్చు.
టైపులకు స్పష్టమైన ``నిర్మాణాత్మక'' అర్థనిర్దేశాలు ఉండటంతో,
అంతఃప్రజ్ఞావాద తర్క సందర్భంలో ఇది మరింత ఆకర్షణీయమైనది. నిజానికి,
ఈ నిర్మాణాత్మక అర్థనిర్దేశాలను చక్కగా స్పష్టం చేసే ఉన్నత-టైపు
అర్థవిచారపు నిర్మాణాత్మక రూపాలను (సాంప్రదాయిక తర్కానికి బదులు
అంతఃప్రజ్ఞావాద తర్కంపై ఆధారపడినవి) అభివృద్ధి చేయవచ్చు; అవి అనేక
విధాలుగా సాంప్రదాయిక ప్రతిరూపాలకంటే ఆసక్తికరమైనవి.

\end{document}
